% Transfer section, version used if endpoint N1 is met.
% First written before any water temperature of the new USGS stations was read; the
% opening paragraph was revised when all results were recomputed with plain integer seeds.
% Numbers are macros filled from
% results/wrr_target_recovery_20261022/e10_sealed_usgs/sealed_results.json.
\subsection{A label-free threshold held the admitted risk at new stations}
The second test used \TNewStations{} USGS stations whose water temperature had not been read by this or the earlier study when the test was designed (\TNewGroups{} independence groups in \TNewHucs{} two-digit hydrologic regions, \TNewYears{} complete station-years). Their eligibility, networks and missingness strata were fixed from the dates of approved daily means alone, and the rule, the models, the interval calibration and the code were frozen before the values were requested. The values below come from the recomputation described in Section~\ref{sec:tests}, made after the first evaluation, so this is no longer a blind test. These are new stations in regions represented in the source, not a new climate or agency.

Without screening, \TNewUnscreenedDeploy\% of D2* values failed under the stations' own gap calendars and \TNewUnscreenedMech\% under long controlled gaps. The frozen rule admits cases in increasing order of the direct failure probability while $k$ times the mean predicted probability among admitted cases stays at or below 5\%, with $k=\TNewFactor$ estimated across source regions; it uses no label of the new stations. It kept the realised risk at \TNewRiskDeploy\% [group-bootstrap upper bound \TNewRiskDeployHigh\%] under real calendars while using \TNewRetDeploy\% of the station-years, and at \TNewRiskMech\% [\TNewRiskMechHigh\%] under long gaps while using \TNewRetMech\%. The preregistered criterion of at most 7\% in both populations was met. Stations with a high share of incomplete years had a realised risk of \TNewRiskDeployHighMissing\% under their own calendars. Across hydrologic regions with at least 200 cases, the worst realised risk was \TNewWorstRegionRisk\%.

The calibrated 90\% intervals covered \TNewHardCoverage\% of the truths where more than three quarters of the climatological peak window was hidden (criterion: at least 85\%; \TNewHardCoverageVerdict). In air-layer cases with the true peak at least half hidden, D2* reduced the cold bias of the plug-in from \TNewPOneBiasPZero{} to \TNewPOneBiasDTwo\,$^{\circ}$C, a reduction of \TNewPOneReduction\% (criterion: at least 50\%; \TNewPOneVerdict). At equal realised risk, the target-aware judgment used \TNewGainMech{} percentage points more cases than the best published rule under long gaps and \TNewGainDeploy{} points more under real calendars. The recovery law predicted a plug-in failure rate of \TNewLawPredicted\% from decision-time information (observed \TNewObservedFail\%).
